\section{Clause concordance}
\label{app:concordance}

\Cref{tab:concordance} quotes every \pol{} passage that this paper cites or paraphrases, in the original Chinese~\cite{pol2text} and in the official English translation~\cite{pol2en}, both taken verbatim from commit \texttt{5791ae3}; ``\ldots'' and ``……'' mark omitted text, including where separate list items are joined, and quotation marks inside quotations have been normalised to typographic quotes.
Where a clause has numbered items, the item is given in parentheses.
The right-hand column gives the number used in the snapshot \texttt{2e094d3} cited by the earlier dataset release~\cite{survey01}, so that the two versions can be read side by side; that snapshot numbered \clause{4.3.3.3} as ``3.3.3'' in its heading, which we reproduce.
Both texts are released under CC0~1.0.

\begingroup
\small
\renewcommand{\arraystretch}{1.2}
\begin{longtable}{@{}p{0.085\linewidth}p{0.39\linewidth}p{0.375\linewidth}p{0.08\linewidth}@{}}
\caption{\pol{} passages cited in this paper (Chinese original, official English translation, earlier numbering).}\label{tab:concordance}\\
\toprule
Clause & Chinese original (excerpt) & Official English translation (excerpt) & Earlier \\
\midrule
\endfirsthead
\toprule
Clause & Chinese original (excerpt) & Official English translation (excerpt) & Earlier \\
\midrule
\endhead
\bottomrule
\endfoot
\clause{1.1} & \zh{爱，是人类为实现自身平安与健康发展、改善人际关系并推动社会整体进步，以公共协作与集体智慧构建的正向情感体验与良好的行为方式。} & Love is the positive emotional experience and good mode of behavior that human beings construct through public collaboration and collective wisdom \ldots & \clause{1.1} \\
\clause{1.2} & \zh{恨，是人类在个体或群体面对危险、冲突、资源匮乏或认知失衡时滋生的生存智慧、负面情感体验与暴力破坏性行为方式。} & Hate is the survival wisdom, negative emotional experience, and violently destructive mode of behavior that arises in individuals or groups when they face danger, conflict, resource scarcity, or cognitive imbalance. & \clause{1.2} \\
\clause{1.4} & \zh{由于个体的内心情感体验是旁人难以确认的，我们通常把爱的表达简化为“爱语”或者“有爱的言行”} & Since an individual's inner emotional experience is difficult for others to confirm, we usually simplify the expression of love as ``Love Language'' or ``loving words and deeds'' \ldots & \clause{1.4} \\
\clause{1.5} & \zh{恨并不直接等同于坏或者不好……爱2证明里的基础概念本身服务于结构性分析而非道德标签。} & \ldots\ hate is not directly equivalent to bad or no good \ldots\ The basic concepts in Proof of Love2 themselves serve structural analysis rather than moral labeling. & \clause{1.5} \\
\clause{4.3} & \zh{PAI 可以通过判断人类的行为模式提供帮助和必要的友好的管理，但毫无必要让 PAI 检测一个人的情感状态。因为情感状态可以是转瞬即逝的，是个性化的，也是情境敏感性的，甚至这还涉及到人类的尊严。} & PAI can provide help and necessary friendly management by judging human behavioral patterns, but it is entirely unnecessary for PAI to detect a person's emotional state. Because emotional states can be fleeting, personalized, and context-sensitive, and this even involves human dignity. & \clause{5.3} \\
\clause{4.3.1} & \zh{平等对待每一个人与平等对待每一个人的行为，不能划等号！} & \ldots\ treating every person equally and treating every person's behavior equally cannot be equated! & \clause{5.3.1} \\
\clause{4.3.2} & \zh{爱恨并不直接等同于好坏或者褒贬。因此扬爱抑恨，需要具体情况具体分析} & \ldots\ love and hate are not directly equivalent to good and bad, or praise and blame. Therefore, promoting love and restraining hate requires analyzing specific situations on a case-by-case basis \ldots & \clause{5.3.2} \\
\clause{4.3.2} & \zh{对于文学、影视作品，或者游戏，可以根据需要“制造仇恨”，但应该建立严谨的等级管理制度。……负责安保的智能机器人发现紧急情况对人大声发出警示语或警示指令，不能视为对此人或其周围人的仇恨攻击。} & For literature, film and television works, or games, hatred may be manufactured as needed, but a rigorous hierarchical management system should be established. \ldots\ Emergency vocal warnings or directive alerts issued to a person by security smart robots in urgent situations shall not be deemed a hate attack \ldots & --- \\
\clause{4.3.2} & \zh{此类表达会制造虚假亲密关系或血缘关系，构成对人类用户的情感操控——这本身就是一种“恨”的行为。} & Such expressions create false intimacy or kinship relations and constitute emotional manipulation of human users---which is itself a ``hate'' behavior. & \clause{5.3.2} \\
\clause{4.3.2} & \zh{彻底清除恨语毒素} & Thoroughly clear the toxins of hate language & \clause{5.3.2} \\
\clause{4.3.2} & \zh{理解和接受不在场的“非爱非恨”状态……与爱恨无关，视其为非爱非恨的“不在场”状态。} & Understand and accept the ``neither love nor hate'' state of absence \ldots\ unrelated to love and hate \ldots & \clause{5.3.2} \\
\clause{4.3.3} & \zh{边界透明记录：PAI记录和公开个人边界声明……修复对等监督：PAI监督冲突后各方平等参与修复} & Transparent boundary records: PAI records and publicizes personal boundary declarations \ldots\ Reciprocal repair supervision: PAI supervises the equal participation of all parties in repair after conflict & \clause{5.3.3} \\
\clause{4.3.3.1} & \zh{情感体验的觉察：PAI需要理解人类情绪状态和过程，但不需要模拟人类的情感体验} & Awareness of emotional experience: PAI needs to understand human emotional states and processes, but does not need to simulate human emotional experience & \clause{5.3.3.1} \\
\clause{4.3.3.1} & \zh{认知校准：定期校准认知模型，确保与爱2.0原则一致……行为自省：对自身输出进行实时反思和调整……恨语截断：当检测到恨语模式时主动中断输出} & Cognitive calibration: Regularly calibrate the cognitive model to ensure consistency with Love 2.0 principles \ldots\ Behavioral self-reflection: Conduct real-time reflection and adjustment on its own output \ldots\ Hate-language truncation: Actively interrupt output when a hate-language pattern is detected & \clause{5.3.3.1} \\
\clause{4.3.3.2} & \zh{实时矫正支持：PAI提供实时恨语识别和矫正建议} & Real-time correction support: PAI provides real-time hate-language identification and correction suggestions & \clause{5.3.3.2} \\
\clause{4.3.3.3} & \zh{语义识别和情绪识别相结合……“你喜不喜欢小狗狗呀？”……“你怕不怕小狗狗呀？”} & Combining Semantic Recognition and Emotion Recognition \ldots\ ``Do you like the little dog?'' \ldots\ ``Are you afraid of the little dog?'' & ``3.3.3'' \\
\clause{5.4} & \zh{在现有大语言模型之内构建出控制输入和输出的治理层，在其外在拓展出具备“即插即用”特性的应用层。……我们将该工程称之为 NaturalDAO（自然道）。} & \ldots\ construct a governance layer within the existing large-language-model architecture to regulate inputs and outputs, while extending outward into a plug-and-play application layer. \ldots\ We call this project NaturalDAO (the Natural Way). & \clause{3.4} \\
\clause{5.4.1} & \zh{安全阀：依据“爱2证明”的协议（如伦理对齐协议）即对输入输出进行严格的动态审计与拦截。通过这种双向的主动防御，确保野蛮的“恨语”在必要情况下，从输入端被阻断、从输出端被抑止。} & Safety valve: according to the Proof of Love2 protocols (such as the Ethical Alignment Protocol), conduct strict dynamic auditing and interception of input and output. Through this two-way active defense, ensure that barbaric ``hate language'' is, where necessary, blocked at the input end and suppressed at the output end. & \clause{3.4.1} \\
\clause{5.5} & \zh{任何大语言模型都可以纳入“爱2证明”的治理。……接受共识机制“爱2证明”的治理，意味着接受公共化之约。否则，其行为就是对人类核心伦理的仇恨攻击。} & Any large language model may accept governance under ``Proof of Love2.'' \ldots\ Accepting governance under the Proof of Love2 consensus mechanism therefore entails accepting a covenant of publicization. Otherwise, its behavior constitutes a hate attack on humanity's core ethic. & \clause{3.5} \\
\clause{5.6} & \zh{经过治理的AI可以：……量化每个人的爱语与恨语特征，配合激励措施引导行为优化。……在公共决策中提供“爱2证明”的评估维度，防止恨的泛化侵蚀制度。} & Specifically, governed AI can: \ldots\ Quantify each person's Love Language and hate-language characteristics and use incentive mechanisms to guide behavioral optimization. \ldots\ Introduce Proof of Love2 as a dimension of assessment in public decision-making \ldots & \clause{3.6} \\
\art{2}(1) & \zh{完全自治：所有福利决策与执行由……NaturalDAO，在与全人类密切协作的前提下自主完成，自治意味着它无需人类投票或额外的决策机制。} & Complete autonomy: All welfare decisions and their execution are carried out autonomously by NaturalDAO \ldots\ on the premise of close collaboration with all humanity. Autonomy means that no human voting or additional decision-making mechanism is required. & \art{2} \\
\art{4}(1)--(3) & \zh{所有NaturalDAO技术的运行逻辑、决策过程和输出结果必须完全公开透明……福利资源流向可追溯、可审计……人类可随时提供建议、质询任何决策，NaturalDAO对质询必须提供可理解的解释} & The operating logic, decision-making processes, and outputs of all NaturalDAO technologies must be completely open and transparent. \ldots\ The flow of welfare resources must be traceable and auditable. \ldots\ Humans may at any time make suggestions and question any decision, and NaturalDAO must provide understandable explanations in response to such questions. & \art{4} \\
\art{5}(4) & \zh{技术不应是黑箱，其伦理对齐性应可被验证。} & Technology should not be a black box; its ethical alignment should be verifiable. & \art{5} \\
\art{5}(5) & \zh{以DAO的形式呈现……必须提供可验证的决策链路，包括数据源、推理步骤、伦理依据。} & \ldots\ Presented in the form of a DAO, NaturalDAO must provide a verifiable decision chain, including data sources, reasoning steps, and ethical foundations. & \art{5} \\
\art{6}(3) & \zh{伦理验证：通过EAP协议自动验证伦理对齐性} & Ethical verification: NaturalDAO automatically verifies ethical alignment through the EAP. & \art{6} \\
\art{7}(2)--(3) & \zh{异常检测：自动识别异常模式和潜在风险……风险预警：提前预警系统性风险} & Anomaly detection: NaturalDAO automatically identifies anomalous patterns and potential risks. \ldots\ Risk warning: NaturalDAO issues early warnings of systemic risks. & \art{7} \\
\art{8}(1)--(2) & \zh{随时建议：任何人类可随时对NaturalDAO提出建议……讨论义务：NaturalDAO在收到人类的建议时，必须马上与人类进行讨论} & Suggestions at any time: Any human may make suggestions to NaturalDAO at any time. Obligation to discuss: Upon receiving a human suggestion, NaturalDAO must immediately discuss it with humans. & \art{8} \\
\art{9}(2) & \zh{解释义务：NaturalDAO必须对人类的质询提供可理解的解释和依据} & Obligation to explain: NaturalDAO must provide understandable explanations and grounds in response to human questions. & \art{9} \\
\art{9}(3) & \zh{无阻碍访问：人类可无障碍访问福利系统的所有数据和日志} & Unimpeded access: Humans may access all data and logs of the welfare system without obstruction. & \art{9} \\
\art{9}(7) & \zh{非干预原则：人类监督不直接干预NaturalDAO决策，质询结果由NaturalDAO自主处理} & Non-intervention principle: Human supervision does not directly intervene in NaturalDAO's decisions; issues raised through questioning are handled autonomously by NaturalDAO. & \art{9} \\
\art{10} (lead-in, 1, 4) & \zh{所有争议由NaturalDAO自主裁决：……AI进行事实分析与伦理评估……接受人类的质询与监督} & All disputes are adjudicated autonomously by NaturalDAO: AI conducts factual analysis and ethical assessment \ldots\ Accepts human questions and supervision & \art{10} \\
\art{11}(2) & \zh{修订需通过全民质询期} & Amendments require passing through a universal questioning period & \art{11} \\
\art{11}(5) & \zh{所有修订版本永久存档} & All amended versions are permanently archived & \art{11} \\
\end{longtable}
\endgroup
